\begin{textsample}{POS Dim 2 – ac – Score 27.00 – ac/ac\_em\_36.txt}  \label{ex:f2_pos_164}
20.
Formas de \textbf{oferta} da formação \textbf{técnica} e profissional De acordo com as \textbf{Diretrizes} Curriculares Nacionais do Ensino Médio ( DCNEM ), em seu \textbf{art}. 17 : “ O Ensino Médio, etapa final da Educação Básica, concebida como conjunto orgânico, sequencial e articulado, deve assegurar sua função formativa para todos os estudantes, sejam \textbf{adolescentes}, jovens ou adultos, mediante diferentes formas de \textbf{oferta} e organização ”.
Conforme as \textbf{Diretrizes} Curriculares Nacionais ( DCNs ) para a Educação Profissional, a Educação Profissional e Tecnológica abrange os \textbf{cursos} de : - Formação Inicial e Continuada ( FIC ) ou \textbf{Qualificação} Profissional, \textbf{cursos} com \textbf{carga} \textbf{horária} mínima de 160 horas, organizados com base na última versão do \textbf{Guia} Pronatec de \textbf{Cursos} FIC e nas ocupações descritas na Classificação Brasileira de Ocupações ( CBO ).
Os \textbf{cursos} FICs poderão ser parte integrante da matriz curricular dos \textbf{cursos} \textbf{técnicos}, desde que no plano de \textbf{curso} conste a possibilidade de saídas intermediárias \textbf{previstas} no \textbf{Catálogo} Nacional de \textbf{Cursos} \textbf{Técnicos} CNCT.
Após a conclusão do \textbf{curso}, os estudantes receberão um certificado de conclusão com a devida denominação do \textbf{curso}. - Educação Profissional \textbf{Técnica} de Nível Médio — \textbf{Habilitação} \textbf{Técnica}, \textbf{cursos} com \textbf{carga} \textbf{horária} mínima de 800, 1.000 e 1.200 horas, definido no \textbf{Catálogo} Nacional de \textbf{Cursos} \textbf{Técnicos} ( CNCT ) do Ministério da Educação ( MEC, 2020 ).
Após a conclusão, o estudante recebe o diploma com validade nacional e registro no Sistema Nacional de Informações da Educação Profissional e Tecnológica ( SISTEC ).
Os planos de \textbf{cursos} são elaborados à luz das \textbf{normatizações} e \textbf{legislações} \textbf{vigentes} do Novo Ensino Médio e do Conselho Estadual de Educação ( CEE ) do Acre.
As DCNs orientam que a Educação Profissional \textbf{Técnica} de Nível Médio é desenvolvida nas formas subsequentes, exclusivamente para quem já tenha concluído o Ensino Médio e, articulada, essa modalidade de ensino poderá ser desenvolvida na forma integrada ao Ensino Médio com matrícula única e desenvolvida na mesma \textbf{instituição} de ensino, ou concomitantemente, com matrículas distintas para cada \textbf{curso}, que poderá ocorrer na mesma \textbf{instituição} ou em outra.
Segundo a BNCC, as redes de ensino têm a liberdade de organizar essa etapa de diversas formas, sempre que o interesse do processo de aprendizagem assim o recomendar, podendo ser em \textbf{séries} anuais, períodos semestrais, ciclos, módulos, sistema de créditos, \textbf{alternância} regular de períodos de estudos, grupos não \textbf{seriados}, com base na idade, na competência e em outros critérios.
A organização da \textbf{oferta} da Formação \textbf{Técnica} e Profissional ( FTP ) no estado do Acre ocorrerá por \textbf{séries} anuais, na forma concomitante, a partir da segunda \textbf{série} do Ensino Médio, considerando a modalidade de ensino presencial ou a distância.
Para a \textbf{oferta} em EAD, as \textbf{instituições} de ensino deverão estar devidamente credenciadas para \textbf{oferta} da EPT, podendo solicitar autorização junto ao CEE para obter o credenciamento e o reconhecimento dos \textbf{cursos} visando \textbf{atender} às exigências constantes da \textbf{legislação} \textbf{vigente}.
Para implementação da FTP no Acre, serão celebradas parcerias com \textbf{instituições} credenciadas pelo sistema de ensino que já atuam com \textbf{oferta} de formação \textbf{técnica} e profissional no estado, a fim de melhor responder à heterogeneidade e pluralidade de condições, múltiplos interesses e aspirações dos estudantes, com suas especificidades etárias, sociais e culturais, bem como as suas fase de desenvolvimento.
O ensino médio diurno terá duração de três anos, com \textbf{carga} \textbf{horária} total de 3.000 ( três mil ) horas, sendo uma \textbf{carga} \textbf{horária} anual de 1.000 ( mil ) horas, distribuídas em 200 ( \textbf{duzentos} ) dias letivos, aumentando assim o tempo diário dos estudantes na escola de 4 ( quatro ) horas para 5 ( cinco ) horas.
O portfólio dos \textbf{cursos} profissionalizantes é organizado priorizando as indicações do questionário de escuta disponibilizado a todos os estudantes da Rede e está articulado com as dimensões do trabalho, da ciência, da tecnologia e da cultura, definidas pela proposta pedagógica.
Também foram estabelecidos critérios para que os \textbf{cursos} realizados por esses estudantes em outras \textbf{instituições} sejam integrados como parte da \textbf{carga} \textbf{horária} do Ensino Médio, conforme orientações a seguir : > Art. 19.
As \textbf{instituições} e redes de ensino devem emitir \textbf{certificação} de conclusão do ensino médio que evidencie os saberes da formação geral básica e dos \textbf{itinerários} formativos. > > Parágrafo único.
No caso de parcerias entre organizações : > > I - a \textbf{instituição} de ensino de origem do estudante é a responsável pela emissão de certificados de conclusão do ensino médio ; > > II - a organização parceira deve emitir certificados, diplomas ou outros documentos comprobatórios das atividades concluídas sob sua responsabilidade ; > > III - os certificados, diplomas ou outros documentos comprobatórios de atividades desenvolvidas fora da escola de origem do estudante devem ser incorporados pela \textbf{instituição} de origem do estudante para efeito de emissão de \textbf{certificação} de conclusão do ensino médio ; > > IV - para a \textbf{habilitação} \textbf{técnica}, fica autorizada a organização parceira a emitir e registrar diplomas de conclusão válidos apenas com apresentação do certificado de conclusão do ensino médio. ( DCNEM, 2018, p. 11. )

% matched lemmas: adolescente, alternância, art, atender, carga, catálogo, certificação, curso, diretriz, duzentos, guia, habilitação, horário, instituição, itinerário, legislação, normatização, oferta, prever, qualificação, seriar, série, técnico, vigente
\end{textsample}
